\section{Auxiliary Discrepancy and Empirical Entropy Bounds}
\label{app:discrepancy}

This appendix collects finite-sample control results that bound the empirical
folded histogram relative to its deterministic limit law. They complement the
structural Parry divergence analysis of the main text by showing that the
sampling error can be made arbitrarily small.

\subsection{Discrepancy of the Rotation Orbit}

The quantitative input is the one-dimensional star discrepancy of the
rotation orbit:
\[
D_N^*(\alpha;x_0)
:=
\sup_{0\le a\le 1}
\left|
\frac{1}{N}\sum_{t=0}^{N-1}\mathbf{1}_{[0,a)}(x_t)-a
\right|.
\]
This quantity is effective because length-$m$ rotation words are determined by a
finite partition of the circle into interval atoms.

\subsection{Total Variation Bounds}

\begin{theorem}[Discrepancy control of microstate and folded laws]\label{thm:tv-certificate-hist}
In the rotation model above, for every $m,N\ge 1$,
\[
\TV(\hat\mu_{m,N},\mu_m)\le 2m\,D_N^*(\alpha;x_0).
\]
Consequently,
\[
\TV(\hat\pi_{m,N},\pi_m^{\mathrm{fold}})
\le
2m\,D_N^*(\alpha;x_0).
\]
If \(\beta=\alpha\), then the sharper bounds
\[
\TV(\hat\mu_{m,N},\mu_m)\le (m+1)D_N^*(\alpha;x_0),
\qquad
\TV(\hat\pi_{m,N},\pi_m^{\mathrm{fold}})\le (m+1)D_N^*(\alpha;x_0)
\]
also hold.
\end{theorem}

\begin{proof}
For fixed $m$, the length-$m$ rotation words are constant on the interval atoms
of the partition of \(\TT\) cut out by the points
\[
0,\ \beta,\ -\alpha,\ \beta-\alpha,\ \dots,\ -(m-1)\alpha,\ \beta-(m-1)\alpha
\pmod 1.
\]
There are at most \(2m\) such endpoints and therefore at most \(2m\) interval
atoms. For any interval \(I=[a,b)\subset\TT\),
\[
\left|
\frac{1}{N}\sum_{t=0}^{N-1}\mathbf{1}_I(x_t)-\leb(I)
\right|
\le
2D_N^*(\alpha;x_0).
\]
Taking half the sum of the atomwise errors gives
\(\TV(\hat\mu_{m,N},\mu_m)\le 2mD_N^*(\alpha;x_0)\).

If \(\beta=\alpha\), the coding is Sturmian and the number of distinct
length-\(m\) factors is exactly \(m+1\)
\cite{MorseHedlund1940,LindMarcus1995}. Summing over the \(m+1\) nonempty
cylinder atoms yields the sharper constant. The folded estimates follow because
deterministic pushforward cannot increase total variation.
\end{proof}

\begin{corollary}[Optimal leading constant within the metallic-mean family]
\label{cor:metallic-mean-golden-unique-optimal-tv-constant}
For $A\ge 1$, let
\[
\alpha_A=[0;A,A,\dots],\qquad
\lambda_A:=\frac{A+\sqrt{A^2+4}}{2},\qquad
\kappa(A):=\frac{A}{\log\lambda_A}.
\]
Then there is a constant $C(A)<\infty$ such that for every $N\ge 2$ and every
$x_0\in\TT$,
\[
D_N^*(\alpha_A;x_0)\le \frac{\kappa(A)\log N+C(A)}{N}.
\]
Consequently, for every fixed resolution $m\ge 1$,
\[
\TV(\hat\pi_{m,N},\pi_m^{\mathrm{fold}})
\le
2m\frac{\kappa(A)\log N+C(A)}{N}.
\]
If \(\beta=\alpha_A\), this sharpens to
\[
\TV(\hat\pi_{m,N},\pi_m^{\mathrm{fold}})
\le
(m+1)\frac{\kappa(A)\log N+C(A)}{N}.
\]
Moreover, $\kappa(A)$ is strictly increasing on $[1,\infty)$, so the golden
branch $A=1$ uniquely minimizes the leading $(\log N)/N$ coefficient among
metallic means.
\end{corollary}

\begin{proof}
Let $p_n/q_n$ be the convergents of $\alpha_A$.  The standard discrepancy bound
\cite{KuipersNiederreiter1974} gives
\[
D_N^*(\alpha_A;x_0)\le \frac{3+nA}{N}
\qquad \text{whenever } q_n\le N<q_{n+1}.
\]
Since $q_{n+1}=Aq_n+q_{n-1}$ and $\lambda_A$ is the Perron root of
$x^2-Ax-1=0$, one has $q_n\ge \lambda_A^{n-1}$ for $n\ge 1$, hence
$n\le 1+\log N/\log\lambda_A$.  Substituting yields the discrepancy estimate
with $C(A)=3+A$.  The TV bound follows from
Theorem~\ref{thm:tv-certificate-hist}.  The monotonicity of
$\kappa(A)=A/\log\lambda_A$ on $[1,\infty)$ is elementary, so its minimum is
attained uniquely at $A=1$.
\end{proof}

\subsection{Cross-Resolution Residuals}

Let $r_{m+1\to m}:\Omega_{m+1}\to\Omega_m$ denote prefix truncation,
$r_{m+1\to m}(u_1\cdots u_{m+1})=u_1\cdots u_m$, and let
$\rho_{m+1\to m}:X_{m+1}\to X_m$ be the analogous truncation on folded words.
Define the projective residual
\[
E_{m,N}
:=
\TV\!\left(\hat\pi_{m,N},(\rho_{m+1\to m})_*\hat\pi_{m+1,N}\right).
\]

\begin{theorem}[Boundary law for the canonical fold]\label{prop:multiscale-residual-by-boundary}
For the canonical fold of Definition~\ref{def:canonical-fold}, let
\[
f_m(\omega):=\Fold_m(r_{m+1\to m}(\omega)),
\qquad
g_m(\omega):=\rho_{m+1\to m}(\Fold_{m+1}(\omega)),
\]
and define the boundary event set
\[
\mathcal{B}_{m+1}:=\{\omega\in\Omega_{m+1}: f_m(\omega)\neq g_m(\omega)\}.
\]
Then
\[
E_{m,N}\le \hat\mu_{m+1,N}(\mathcal{B}_{m+1}).
\]
Moreover, if \(\omega=u_1\cdots u_{m+1}\) and \(v:=u_1\cdots u_m\), then
\[
\omega\in\mathcal{B}_{m+1}
\iff
u_{m+1}=1\ \text{and}\ N(v)\ge F_{m+1}.
\]
\end{theorem}

\begin{proof}
Both $\hat\pi_{m,N}$ and $(\rho_{m+1\to m})_*\hat\pi_{m+1,N}$ are pushforwards
of $\hat\mu_{m+1,N}$, respectively by $f_m$ and $g_m$.  Therefore
\[
\TV((f_m)_*\nu,(g_m)_*\nu)\le \nu(f_m\neq g_m)
\]
for any probability measure $\nu$, and the first claim follows with
$\nu=\hat\mu_{m+1,N}$.

For the second claim, write \(N:=N(v)\). Since
\[
0\le N\le \sum_{k=1}^m F_{k+1}=F_{m+3}-2,
\]
the Zeckendorf expansion of \(N\) has the form
\[
Z(N)=(c_1,\dots,c_{m+1},0,0,\dots).
\]
If \(u_{m+1}=0\), then \(N(\omega)=N\), so \(f_m(\omega)=g_m(\omega)\).

Assume next that \(u_{m+1}=1\), so \(N(\omega)=N+F_{m+2}\). If \(N<F_{m+1}\),
then \(c_m=c_{m+1}=0\), hence
\[
Z(N+F_{m+2})=(c_1,\dots,c_m,1,0,0,\dots)
\]
is already a legal Zeckendorf expansion. Therefore
\[
g_m(\omega)=\rho_{m+1\to m}(\Fold_{m+1}(\omega))=(c_1,\dots,c_m)=f_m(\omega),
\]
so \(\omega\notin\mathcal{B}_{m+1}\).

Conversely, suppose \(u_{m+1}=1\) and \(\omega\notin\mathcal{B}_{m+1}\). Write
\[
Z(N+F_{m+2})=(d_1,\dots,d_{m+2},0,0,\dots).
\]
Since the first \(m\) digits agree, one has \(d_j=c_j\) for \(1\le j\le m\).
Subtracting the two Zeckendorf expansions gives
\[
F_{m+2}
=
(d_{m+1}-c_{m+1})F_{m+2}+d_{m+2}F_{m+3}.
\]
Because \(F_{m+3}>F_{m+2}\), necessarily \(d_{m+2}=0\) and
\(d_{m+1}-c_{m+1}=1\). Hence \(c_{m+1}=0\) and \(d_{m+1}=1\). Since
\((d_1,\dots,d_{m+2})\) is a legal Zeckendorf digit string and \(d_m=c_m\), one
must also have \(c_m=d_m=0\). Thus \(Z(N)\) has no nonzero digit in positions
\(m\) or \(m+1\), which is equivalent to \(N<F_{m+1}\).

We have proved
\[
\omega\notin\mathcal{B}_{m+1}
\iff
u_{m+1}=0\ \text{or}\ \bigl(u_{m+1}=1\ \text{and}\ N(v)<F_{m+1}\bigr),
\]
which is exactly the stated characterization.
\end{proof}

\begin{corollary}[Decomposition relative to the Parry measure]\label{cor:tv-decomposition-parry}
For every $m,N\ge 1$,
\[
\TV(\hat\pi_{m,N},\pi_m)
\le
2m\,D_N^*(\alpha;x_0)+\TV(\pi_m^{\mathrm{fold}},\pi_m).
\]
If \(\beta=\alpha\), one may replace \(2m\) by \(m+1\).
\end{corollary}

\begin{proof}
Theorem~\ref{thm:tv-certificate-hist} gives
$\TV(\hat\pi_{m,N},\pi_m^{\mathrm{fold}})\le 2mD_N^*(\alpha;x_0)$, and the
triangle inequality completes the proof.
\end{proof}

\subsection{Relative-Entropy Bounds}

\begin{lemma}[Relative entropy from total variation under strict positivity]\label{lem:kl-from-tv-qmin}
Let $p,q$ be probability measures on a finite set $S$, assume $p\ll q$, and
suppose that $q$ is strictly positive on $S$. Let
\[
q_{\min}:=\min_{s\in S}q(s)>0.
\]
Then
\[
\KL(p\|q)
\le
\log\!\left(1+\frac{\|p-q\|_1^2}{q_{\min}}\right)
\le
\frac{1}{q_{\min}}\|p-q\|_1^2
=
\frac{4}{q_{\min}}\TV(p,q)^2.
\]
\end{lemma}

\begin{proof}
The standard inequality \(\KL\le \log(1+\chi^2)\) gives
\[
\KL(p\|q)\le \log\bigl(1+\chi^2(p,q)\bigr),
\qquad
\chi^2(p,q):=\sum_{s\in S}\frac{(p(s)-q(s))^2}{q(s)}.
\]
Since \(q(s)\ge q_{\min}\) for every \(s\in S\),
\[
\chi^2(p,q)
\le
\frac{1}{q_{\min}}\sum_{s\in S}(p(s)-q(s))^2
\le
\frac{1}{q_{\min}}\|p-q\|_1^2.
\]
The second inequality uses \(\log(1+u)\le u\). See
\cite{SasonVerdu2016fDivergence}.
\end{proof}

\begin{definition}[Diophantine lower bound]
For a badly approximable irrational $\alpha$, let $c(\alpha)>0$ be a constant
such that
\[
\|k\alpha\|\ge \frac{c(\alpha)}{k}\qquad (k\ge 1),
\]
where $\|x\|:=\min_{n\in\mathbb{Z}}|x-n|$.
\end{definition}

\begin{proposition}[Lower bound on the smallest microstate atom]\label{prop:qmin-lowerbound-markov}
Assume $\beta=\alpha$ and let
\[
q_{\min}(m):=\min\{\mu_m(u):\mu_m(u)>0\}.
\]
Then
\[
q_{\min}(m)\ge \min_{1\le k\le m}\|k\alpha\|\ge \frac{c(\alpha)}{m}.
\]
For the golden slope $\alpha=\varphi^{-1}$ one may take
$c(\alpha)=1/\sqrt5$, hence
\[
q_{\min}(m)\ge \frac{1}{\sqrt5\,m}.
\]
\end{proposition}

\begin{proof}
When $\beta=\alpha$, the partition endpoints for length-$m$ microstates reduce
to the set $\{-k\alpha:0\le k\le m\}$.  Thus the positive atoms of $\mu_m$ are
exactly the adjacent gaps determined by these points, and the smallest one is
at least $\min_{1\le k\le m}\|k\alpha\|$.  The badly approximable lower bound
then yields $c(\alpha)/m$.  For $\varphi^{-1}$ the optimal Markov constant is
$1/\sqrt5$.
\end{proof}

\begin{corollary}[Microstate relative-entropy bound]\label{cor:rotation-microstate-kl-certificate}
Assume $\beta=\alpha$ and $\alpha$ is badly approximable.  Then
\[
\KL(\hat\mu_{m,N}\|\mu_m)
\le
\frac{4}{q_{\min}(m)}\TV(\hat\mu_{m,N},\mu_m)^2
\le
\frac{4(m+1)^2}{q_{\min}(m)}D_N^*(\alpha;x_0)^2.
\]
For $\alpha=\varphi^{-1}$ this becomes
\[
\KL(\hat\mu_{m,N}\|\mu_m)
\le
4\sqrt5\,m(m+1)^2 D_N^*(\alpha;x_0)^2.
\]
\end{corollary}

\begin{proof}
Let \(S_m:=\supp(\mu_m)\). In the Sturmian slice \(\beta=\alpha\), every
observed length-\(m\) word belongs to the Sturmian language, hence
\(\hat\mu_{m,N}\) and \(\mu_m\) may both be viewed as probability measures on
\(S_m\). On this set, \(\mu_m\) is strictly positive and
\[
\min_{u\in S_m}\mu_m(u)=q_{\min}(m).
\]
The claim therefore follows from Theorem~\ref{thm:tv-certificate-hist},
Lemma~\ref{lem:kl-from-tv-qmin}, and
Proposition~\ref{prop:qmin-lowerbound-markov}.
\end{proof}

\begin{lemma}[KL monotonicity under pushforward]\label{lem:kl-monotone-pushforward}
Let $F:S\to T$ be a map between finite sets and let $p,q$ be probability
measures on $S$ with $p\ll q$.  Then
\[
\KL(F_*p\|F_*q)\le \KL(p\|q).
\]
\end{lemma}

\begin{proof}
This is the data-processing inequality for relative entropy under deterministic
coarse graining.
\end{proof}

\begin{corollary}[Folded relative-entropy bound]\label{cor:rotation-folded-kl-certificate}
Under the assumptions of
Corollary~\ref{cor:rotation-microstate-kl-certificate},
\[
\KL(\hat\pi_{m,N}\|\pi_m^{\mathrm{fold}})
\le
\KL(\hat\mu_{m,N}\|\mu_m)
\le
\frac{4(m+1)^2}{q_{\min}(m)}D_N^*(\alpha;x_0)^2.
\]
\end{corollary}

\begin{proof}
Apply Lemma~\ref{lem:kl-monotone-pushforward} with $F=\Fold_m$.
\end{proof}
